\subsection{Hopcroft-Karp: matching bipartito}\label{rec:hopcroft}

Matching máximo en un grafo bipartito. Izquierda: \codigo{0..n-1}; derecha:
\codigo{0..m-1}. El cover usa \codigo{0..n-1} para la izquierda y
\codigo{n..n+m-1} para la derecha.

\tablafunc{
\codigo{HK(n, m)} & crea las dos particiones & $O(n+m)$ \\
\codigo{add(u, v)} & arista izquierda \codigo{u} -- derecha \codigo{v} & $O(1)$ \\
\codigo{matching()} & tamaño del matching máximo & $O(E\sqrt V)$ \\
\codigo{cover()} & cover mínimo (Kőnig) & $O(V+E)$ \\
\codigo{indep()} & conjunto independiente máximo & $O(V+E)$ \\
\codigo{path_cover} & cobertura mínima de caminos en un DAG & $O(E\sqrt V)$ \\
}

\copiar{el bloque \texttt{COPIA HOPCROFT-KARP}.}
\lstinputlisting{codigo/hopcroft_karp.cpp}

\ojo{Si solo necesitas el tamaño del matching, \codigo{matching()} es
suficiente. El cover y el independiente son complementarios: su tamaño suma
\codigo{n + m}.}